\section{Ambiguity and robust choice: formal proofs}
\label{sec:annex-p17-ambiguity-and-robust-choice}

\begin{annexitem}{annex:p17.ambiguity.seu.benchmark}{thm:p17-ambiguity-seu-benchmark}{Pasting, common prizes and the probability comparisons}
\begin{proof}
Finite summation gives
\[
 E_p(f_Eh)-E_p(g_Eh)
                     =\sum_{s\in E}p(s)[x_f(s)-x_g(s)].
\]
The common outside act \(h\) has cancelled.
The right side is unchanged if \(h\) is replaced by \(h'\),
so its sign gives the sure-thing equivalence for weak,
strict and indifferent comparisons. The same identity
covers the empty event, the whole state space, and
events of probability zero.

For the binary bet, split the finite sum into \(E\)
and its complement:
\[
 E_p(b_E)=p(E)u(H)+(1-p(E))u(L)
                      =u(L)+(u(H)-u(L))p(E).
\]
The positive coefficient \(u(H)-u(L)\) makes the
ordering exactly that of \(p(E)\).
Since \(p(S\setminus E)=1-p(E)\), complementation
reverses each strict comparison and preserves equality.
If \(D\) is disjoint from both \(E,F\), finite additivity
gives \(p(E\cup D)-p(F\cup D)=p(E)-p(F)\),
proving the common-disjoint-event assertion.

To see the final qualification directly, use two states,
utilities zero and one, and rank utility vectors by their
minimum coordinate. The complementary bets \((1,0)\)
and \((0,1)\) are indifferent.
But \((1,1)\) is strictly preferred to \((0,1)\);
changing their common second coordinate from one to
zero makes \((1,0)\) and \((0,0)\) indifferent.
This preference violates the sure-thing condition
just proved necessary for every one-prior model.
Thus the one observed indifference did not establish
the representation of the whole preference.
\end{proof}
\end{annexitem}

\begin{annexitem}{annex:p17.ellsberg.two.urns}{thm:p17-ellsberg-two-urns}{All orientations of the two-urn contradiction}
\begin{proof}
Suppose one additive subjective assessment represents
the four bets, and write \(q_{\rm I},q_{\rm II}\)
for its probabilities of red in the two urns.
The black probabilities must be \(1-q_{\rm I}\)
and \(1-q_{\rm II}\), since each urn has only
these two colors.
By the binary-bet formula, preference for known-urn
red requires \(q_{\rm II}>q_{\rm I}\).
Preference for known-urn black requires
\(1-q_{\rm II}>1-q_{\rm I}\), or
\(q_{\rm II}<q_{\rm I}\). These are incompatible.
Swapping the two urns proves the reversed case.

With weak preferences, the inequalities become
\(q_{\rm II}\geq q_{\rm I}\) and
\(q_{\rm II}\leq q_{\rm I}\), so equality is forced;
either strict preference would then be impossible.
This argument did not require the within-urn
indifferences, or even calibration of \(q_{\rm II}\)
to the known physical fraction. If all bets are
indifferent, take \(q_{\rm I}=q_{\rm II}=1/2\).
For example, the joint prior assigning one quarter
to each auxiliary pair of colors realizes these
marginals and all the indifferences, while agreeing
with the known urn information.
An unknown composition has not thereby been proved
to be half red; this is a possible subjective
representation of that particular pattern.
\end{proof}
\end{annexitem}

\begin{annexitem}{annex:p17.ellsberg.marked.arrow}{thm:p17-ellsberg-marked-arrow}{Every substitution in the marked and Arrow arguments}
\begin{proof}
In a transitive weak preference, replacing an act by
an indifferent act preserves a weak comparison:
if \(x\sim y\), transitivity of \(x\succeq y\succeq z\)
and \(y\succeq x\succeq z\) gives
\(x\succeq z\) if and only if \(y\succeq z\).
The same statement applies to the second argument,
and to strict comparisons by using both weak directions.

In the marked experiment, put \(K=(H,L,L,H)\).
The sure-thing condition permits changing the common
fourth consequence of \(M_1,M_3\) from \(L\) to \(H\),
giving
\[
       M_1\succeq M_3\quad\Longleftrightarrow\quad K\succeq M_6.
\]
Since \(M_6\sim M_5\), the right side is equivalent
to \(K\succeq M_5\).
These latter acts share first consequence \(H\).
Changing that common consequence to \(L\) gives
\[
                    K\succeq M_5
                         \quad\Longleftrightarrow\quad M_4\succeq M_2.
\]
By \(M_4\sim M_3\) and \(M_2\sim M_1\), this last
comparison is equivalent to \(M_3\succeq M_1\).
Consequently
\[
                   M_1\succeq M_3
                           \quad\Longleftrightarrow\quad M_3\succeq M_1.
\]
Completeness supplies at least one of these weak
comparisons, so the equivalence supplies both.
Thus \(M_1\sim M_3\), and the two original
indifferences give all four in one indifference class.
This proves the entire marked-ball argument, with
the common columns and both indifferent substitutions
identified.

For Arrow's variant, \(A_1\sim A_4\),
\(A_2\sim A_3\), and \(A_1\succ A_2\)
give \(A_4\succ A_3\) by those substitutions.
The compared acts \(A_1,A_2\) agree at
\(R_{\rm I}\), where both pay \(H\), and at
\(B_{\rm II}\), where both pay \(L\).
Change these common entries to \(L,H\), respectively,
leaving the two middle entries unchanged.
The resulting pair is exactly \(A_3,A_4\).
The sure-thing equivalence, in both weak directions,
preserves strict preference, so yields \(A_3\succ A_4\).
This contradicts the strict comparison already forced
by transitivity. No numerical probabilities or
utility arithmetic were used in either proof.
\end{proof}
\end{annexitem}

\begin{annexitem}{annex:p17.ellsberg.three.colors}{thm:p17-ellsberg-three-colors}{The common yellow column and additive impossibility}
\begin{proof}
Take \(E=\{R,B\}\). The acts \(F_1,F_2\) have
common consequence \(L\) on the outside event \(\{Y\}\).
Replacing that common consequence by \(H\) gives
\(F_3,F_4\). Thus the sure-thing condition would
preserve the strict ordering of the first pair in
the second pair. Both source patterns reverse it,
so both violate that condition.

For the explicit subjective expected-utility
contradiction, put \(\Delta=u(H)-u(L)>0\).
Then
\[
\begin{aligned}
 E_p(F_1)-E_p(F_2)&=\Delta[p(R)-p(B)],\\
 E_p(F_3)-E_p(F_4)
       &=\Delta[p(R)+p(Y)-p(B)-p(Y)]
         =\Delta[p(R)-p(B)].
\end{aligned}
\]
The same real number cannot have opposite strict
signs. In particular the principal pattern would
require both \(p(R)>p(B)\) and \(p(R)<p(B)\).
This proof does not use \(p(R)=1/3\), so assigning
some other subjective red probability cannot
repair the representation.
It uses only the ordering of the two consequence
utilities, not their numerical spacing or the
curvature of utility at other monetary outcomes.
The conclusion excludes exact representation of
the stated comparisons. No distance on noisy
choice rules or approximate representations was
specified, so no quantitative approximation
claim follows from this argument.
\end{proof}
\end{annexitem}

\begin{annexitem}{annex:p17.finite.maxmin.properties}{thm:p17-finite-maxmin-properties}{Every property of the specified finite functional}
\begin{proof}
Each \(E_{p^j}(f)\) is a finite real sum.
A nonempty finite set of such numbers has a
minimum. The values on a nonempty finite menu
likewise have a maximum. Comparison of these
real values is complete and transitive.
Every expectation is a convex combination of
the state utilities, so lies between their
minimum and maximum. The same bounds hold
after taking the minimum over priors.

If \(x_f(s)\geq x_g(s)\) at every state,
nonnegative prior weights give
\(E_p(f)\geq E_p(g)\) for every \(p\).
Taking minima proves monotonicity.
For \(d=\|x_f-x_g\|_\infty\), each prior gives
\[
             |E_p(f)-E_p(g)|
                     \leq\sum_sp(s)d=d.
\]
Thus \(E_p(f)\geq E_p(g)-d\geq V_{\mathcal C}(g)-d\).
Taking its minimum and then interchanging \(f,g\)
proves the Lipschitz inequality.

If \(q=\sum_j\lambda_jp^j\) is in the convex hull,
\(\lambda_j\geq0,\sum_j\lambda_j=1\), then
\[
               E_q(f)=\sum_j\lambda_jE_{p^j}(f)
                                      \geq\min_j E_{p^j}(f).
\]
A listed minimizing prior belongs to the hull,
so the minima are equal and attained.
The hull is convex by combining its coefficient
vectors. To check its closedness explicitly,
let \(q_n\) in the hull tend to \(q\), and
choose coefficient vectors \(\lambda^{(n)}\)
in the finite simplex. A bounded sequence in
\([0,1]\) has a convergent subsequence: repeatedly
select a closed half interval containing infinitely
many terms, with nested interval lengths tending
to zero, and take increasingly indexed terms
from them. Apply this construction successively
to the finitely many coefficient coordinates.
A subsequence of coefficient vectors then
converges to \(\lambda\), with every
\(\lambda_j\geq0\) and \(\sum_j\lambda_j=1\).
Finite summation gives \(q=\sum_j\lambda_jp^j\),
so \(q\) remains in the hull. A point in the
Euclidean closure is the limit of points chosen
within distance \(1/n\), proving closedness.
This closed finite convex hull is the asserted
polytope.
Every such prior also gives a probability measure on \(2^S\):
any disjoint countable family of subsets has at most
\(|S|\) nonempty members, so countable additivity reduces
to the already specified finite sums.
The source's affine lottery utility is exactly
\(\widehat u(y)=\sum_z y(z)u(z)\), and its integral of
\(\widehat u\circ f\) is the displayed \(E_p(f)\).

If \(\mathcal C\subseteq\mathcal D\) are nonempty finite
prior sets, taking a
minimum over the larger set cannot increase
an act's value. Rankings may nevertheless change:
with vectors \(f=(1,1)\), \(g=(0,3)\), the set
\(\{\delta_2\}\) prefers \(g\), whereas
\(\{\delta_1,\delta_2\}\) prefers \(f\).
This example is realized by consequences with
utilities \(0,1,3\).
Finally replacing \(u\) by \(au+b\), \(a>0\),
replaces \(x_f\) by \(ax_f+b\mathbf1\).
Each prior expectation becomes \(aE_p(f)+b\);
taking the minimum yields \(aV_{\mathcal C}(f)+b\).
If \(u\) is constant all expectations are that
constant. If it is nonconstant, two constant
acts at consequences with different utility
are strictly ranked under every prior and
under their minimum.
\end{proof}
\end{annexitem}

\begin{annexitem}{annex:p17.robust.dominance}{thm:p17-robust-dominance}{Robust and worst-case comparisons with all strictness cases}
\begin{proof}
The inequalities defining robust dominance are
reflexive, and two such inequalities under
each fixed prior compose transitively.
If they hold for \(f,g\), then for every prior
\(E_p(f)\geq E_p(g)\geq V_{\mathcal C}(g)\).
Taking the minimum gives
\(V_{\mathcal C}(f)\geq V_{\mathcal C}(g)\).
If every priorwise comparison is strict,
the finite set of positive differences has
a positive minimum
\[
               \delta=\min_{p\in\mathcal C}[E_p(f)-E_p(g)]>0.
\]
The same argument yields
\(V_{\mathcal C}(f)\geq V_{\mathcal C}(g)+\delta\).
Finiteness and strictness under every retained
prior justify this positive common margin.

For the examples take two states with priors
\(\delta_1,\delta_2\), and consequences allowing
the displayed utility values. The expectations
of \((1,0)\) under these priors are \(1,0\);
those of \((0,1)\) are \(0,1\).
Each beats the other under one prior, so
neither robustly dominates.
The vector \((1,1)\) has worst-case value
one, while \((0,2)\) has value zero.
Under \(\delta_2\), however, their values
are one and two, disproving the converse
from maxmin preference to robust dominance.
Finally \((0,2)\) weakly dominates \((0,1)\)
under both priors and strictly under
\(\delta_2\), but both minimum values are
zero. This proves the qualification on
strictness under only one prior.
\end{proof}
\end{annexitem}

\begin{annexitem}{annex:p17.maxmin.mixtures}{thm:p17-maxmin-mixtures}{The exact forward mixture properties}
\begin{proof}
The finite lottery-mixture formula gives
\[
 x_{\lambda f+(1-\lambda)g}
                       =\lambda x_f+(1-\lambda)x_g.
\]
Consequently, for every prior \(p\),
\[
 E_p(\lambda f+(1-\lambda)g)
   =\lambda E_p(f)+(1-\lambda)E_p(g)
   \geq\lambda V_{\mathcal C}(f)+(1-\lambda)V_{\mathcal C}(g).
\]
Taking the minimum proves concavity, including
the endpoints \(\lambda=0,1\).
For a constant act \(h\), every prior gives
the same \(E_p(h)=x_h\). Its term therefore
factors out of the minimum, proving the
displayed equality.
For \(0<\lambda<1\), differences between
the two constant mixtures are exactly
\(\lambda[V_{\mathcal C}(f)-V_{\mathcal C}(g)]\).
Their signs prove certainty independence
for weak, strict and indifferent comparisons.
If \(f\sim g\), concavity gives
\(V_{\mathcal C}(\lambda f+(1-\lambda)g)
\geq V_{\mathcal C}(f)\), which is the
source's uncertainty-aversion property.

For mixture continuity, suppose
\(V(f)>V(g)>V(h)\), abbreviating \(V_{\mathcal C}\)
by \(V\), and put \(D=\|x_f-x_h\|_\infty\).
Here \(D>0\), since otherwise \(V(f)=V(h)\).
Let \(d_1=V(f)-V(g)>0\) and
\(d_2=V(g)-V(h)>0\). Set
\[
 \alpha=1-\min\{1/2,d_1/(2D)\},\qquad
 \beta=\min\{1/2,d_2/(2D)\}.
\]
Both numbers are strictly between zero and
one. The Lipschitz bound gives
\[
\begin{aligned}
 V(\alpha f+(1-\alpha)h)&\geq V(f)-(1-\alpha)D
                                         \geq V(f)-d_1/2>V(g),\\
 V(\beta f+(1-\beta)h)&\leq V(h)+\beta D
                                         \leq V(h)+d_2/2<V(g).
\end{aligned}
\]
This proves exactly the mixture-continuity
condition stated in the body.
All these calculations start with a specified
utility and prior set. They establish the
forward properties and contain no converse
representation or uniqueness argument.
\end{proof}
\end{annexitem}

\begin{annexitem}{annex:p17.maxmin.full.independence}{thm:p17-maxmin-full-independence}{The complete single-prior characterization for independence}
\begin{proof}
If every listed prior equals \(p\), the
functional is \(E_p\), which is affine under
pointwise lottery mixtures. Thus
\[
 E_p(\lambda f+(1-\lambda)h)
       -E_p(\lambda g+(1-\lambda)h)
                        =\lambda[E_p(f)-E_p(g)].
\]
For \(\lambda>0\) this proves full independence.

Conversely suppose there are two distinct
priors. Some state \(s_0\) has different
probabilities under them. Let
\[
       l=\min_{p\in\mathcal C}p(s_0),\qquad
       q=\max_{p\in\mathcal C}p(s_0),
                           \quad 0\leq l<q\leq1.
\]
Since \(u\) is nonconstant, choose two
consequences with unequal utility and apply
a positive affine normalization making their
utilities zero and one. This leaves the
preference unchanged by the previous theorem.
Use only these two consequences in the
following acts.
Let \(f\) give utility \(\mathbf1_{\{s_0\}}\),
let \(h\) give its complement
\(1-\mathbf1_{\{s_0\}}\), and let \(g\)
be the constant lottery giving the high
consequence with probability \(l\).
Then \(V(f)=l=V(g)\).
The half mixture of \(f,h\) is the
constant lottery of utility \(1/2\), so
has value \(1/2\).
The half mixture of \(g,h\) has value
\[
          \frac l2+\frac12
                   \min_{p\in\mathcal C}[1-p(s_0)]
                       =\frac{1+l-q}{2}<\frac12.
\]
Full independence would preserve the
indifference \(f\sim g\) after the
same half mixture with \(h\), a
contradiction. The entire lottery-act
domain supplies the constant lottery
with high-outcome probability \(l\);
this step need not be available in
an arbitrary restricted menu.

For the stated two-state example,
the minimum coordinate of each
complementary bet is zero, and their
half mixture has both coordinates
\(1/2\). With \(h=f\), independence
would preserve \(f\sim g\) when
comparing \(f\) to \((f+g)/2\),
but their values are zero and one half.
If \(u\) is constant the premise used
to select distinct consequence utilities
fails and all mixtures are indifferent,
as explicitly excluded from the iff
statement.
\end{proof}
\end{annexitem}

\begin{annexitem}{annex:p17.maxmin.urns}{thm:p17-maxmin-urns}{All prior tables and the distinct compatible sets}
\begin{proof}
For the three-color family write
\(p_\theta=(1/3,\theta,2/3-\theta)\),
\(0\leq\theta\leq2/3\).
The four normalized expectations are
\[
             1/3,\qquad \theta,\qquad 1-\theta,\qquad 2/3.
\]
Their minima over the two endpoints
are \(1/3,0,1/3,2/3\).
Every intermediate expectation of any
act is affine in \(\theta\), and
therefore lies between its two endpoint
expectations. Thus adding the 59 interior
physical compositions or the whole
convex interval changes no minimum.
This proves the precise information
bridge while imposing no probability
distribution over the possible
compositions.

For the auxiliary two-urn space,
a joint distribution with known-urn
red marginal \(1/2\) necessarily has
the form
\[
                    (a,b,1/2-a,1/2-b),
                              \qquad 0\leq a,b\leq1/2.
\]
Put \(\xi=2a,\zeta=2b\).
It is a convex combination of the
four listed corner distributions:
give weights \(\xi\zeta\) to the
corner \((a,b)=(1/2,1/2)\),
\((1-\xi)(1-\zeta)\) to \((0,0)\),
\(\xi(1-\zeta)\) to \((1/2,0)\),
and \((1-\xi)\zeta\) to \((0,1/2)\).
These weights are nonnegative and
sum to one, and their two coordinates
average to \(a,b\). Conversely every
listed corner has the required known
marginal, which is preserved by
convex combinations. This proves the
claimed exact convex hull.

The unknown-red event has probability
\(a+b\), ranging from zero to one,
and unknown black has its complement.
Each therefore has minimum zero.
Known red and known black both have
probability \(1/2\) at every prior,
giving the other two values.
The auxiliary coupling is a device
for comparing the four event bets;
it asserts no physical correlation
between two actual draws.
If independent uniform draws from
the two urns are additionally
specified, the same four marginal-bet
values result: the extreme unknown
compositions already attain the
ambiguous minima, and the known
marginal stays one half.

For the marked urn, an endpoint
prior has masses \((\theta,1/2-\theta,
1/4,1/4)\), with \(\theta=0\) or \(1/2\).
Reading the six event indicators
from the table gives expectations
\[
          \theta,\ 1/2-\theta,\ 1/4,\ 1/4,\ 1/2,\ 1/2.
\]
Taking minima proves the six values.
Arrow's four indicators instead give
\[
             1/2,\ \theta+1/4,\ 3/4-\theta,\ 1/2,
\]
with respective minima
\(1/2,1/4,1/4,1/2\). Thus the
specified maxmin preference satisfies
the displayed indifferences and
strict comparisons, while the
preceding pasting proof identifies
the failed axiom.

Finally the two \(d\)-dependent
priors in the statement are
probability vectors because
\(0<d\leq1/3\).
In the four expectation formulas,
their black probabilities are
\(1/3-d\) and \(1/3+d\).
Taking their minima gives
\((1/3,1/3-d,2/3-d,2/3)\).
Both relevant gaps equal \(d>0\).
Different \(d\)'s give different
prior sets and value gaps while
preserving all those strict
qualitative comparisons. Hence
those comparisons cannot identify
one of these sets.
\end{proof}
\end{annexitem}

\begin{annexitem}{annex:p17.ambiguity.randomization.averaging}{thm:p17-ambiguity-randomization-averaging}{Both finite mixture calculations}
\begin{proof}
Let \(q\) be the red probability
under any fixed unknown-urn composition.
A fair coin independent of the draw
chooses the red bet with probability
one half and the black bet with
probability one half. The winning
probability is
\[
                        \tfrac12 q+\tfrac12(1-q)=\tfrac12.
\]
Equivalently, at each possible color
state the resulting act pays a
half-high, half-low lottery. Thus
it is a constant lottery act on
the normalized utility scale and
has expectation, and maxmin value,
one half under every prior.
Its availability changes the menu:
the original comparisons between
specified pure-color bets neither
offer this act nor assert that it
would be rejected.

For fixed weights over priors,
finite rearrangement gives
\[
 \sum_jw_jE_{p^j}(f)
       =\sum_jw_j\sum_sp^j(s)x_f(s)
       =\sum_s\left(\sum_jw_jp^j(s)\right)x_f(s).
\]
The vector in parentheses is
nonnegative and sums to one, so
is exactly the prior \(\bar p\).
The same vector is used for every
act because the weights were fixed.
The single-prior benchmark therefore
proves the sure-thing condition
and excludes the strict source
reversal. A different minimizing
prior for different acts is not
this fixed averaging operation.
\end{proof}
\end{annexitem}

\begin{annexitem}{annex:p17.ellsberg.confidence.rule}{thm:p17-ellsberg-confidence-rule}{The original numerical indices and every endpoint}
\begin{proof}
For a fixed act,
\[
\begin{split}
 W_\rho(f)
 &=\min_{p\in\mathcal C}
             [\,\rho E_{p^0}(f)+(1-\rho)E_p(f)\,]\\
 &=\min_{p\in\mathcal C}
                    E_{\rho p^0+(1-\rho)p}(f).
\end{split}
\]
The first equality uses \(1-\rho\geq0\);
the second uses linearity in the prior.
Every vector \(\rho p^0+(1-\rho)p\)
is a prior. This proves the finite
maxmin bridge, including \(\rho=0\)
and \(\rho=1\).

On utility scale zero and six,
the four estimated expectations
under the symmetric \(p^0\) are
\((2,2,4,4)\). The endpoint minima
are \((2,0,2,4)\), six times
the previously calculated values.
Their weighted average is exactly
\[
                         (2,2\rho,2+2\rho,4).
\]
Substitution gives the two original
tables at \(1/4\) and \(3/4\).
The gaps favoring \(F_1\) over
\(F_2\) and \(F_4\) over \(F_3\)
are both \(2(1-\rho)\), positive
for every \(\rho<1\) and zero
at \(\rho=1\).

For \(F_2\), the worst endpoint
is \((1/3,0,2/3)\); for \(F_3\)
it is \((1/3,2/3,0)\).
Mixing each with \(p^0\) gives
the two displayed effective priors.
At \(\rho=1/4\) these are
\((1/3,1/12,7/12)\) and
\((1/3,7/12,1/12)\).
At \(\rho=3/4\) they are
\((1/3,1/4,5/12)\) and
\((1/3,5/12,1/4)\).
These computations reproduce the
source's later weight tables
after the explicit exchange of
its two unknown-color names.
The set \(\mathcal C_\rho\)
is fixed across all the acts;
only the member attaining a
particular minimum varies.
Both \(E_{p^0}(f)\) and \(V_{\mathcal C}(f)\)
lie between \(\min_z u(z)\) and \(\max_z u(z)\).
The identity
\[
 W_\rho(f)-E_{p^0}(f)
        =(1-\rho)[V_{\mathcal C}(f)-E_{p^0}(f)]
\]
therefore gives the asserted uniform bound.
The two already calculated strict gaps are
\(2(1-\rho)>0\) at every \(\rho<1\), so
this convergence of values coexists with
the exact qualitative reversals.

For the optimistic variant, the
endpoint maxima are \((2,4,6,4)\).
Mixing with estimated values
\((2,2,4,4)\) gives
\((2,4-2\rho,6-2\rho,4)\).
The two gaps for the reversed
pattern are again \(2(1-\rho)\).
This is an explicit calculation
of the source's mentioned variant,
not a claim that maxmin and
maximax are the same criterion.

In the two-urn tradeoff, every
reasonable prior assigns the
known-red bet probability \(q\),
so both its estimated expectation
and security level are \(q\).
For unknown red the stipulated
best estimate is \(1/2\) and
the security level is zero.
Their normalized weighted indices
are therefore \(q\) and \(\rho/2\).
Strict preference, equality, and
reverse preference occur according
as \(q\) is greater than, equal to,
or less than \(\rho/2\).
The example \(q=49/100\) from the
source can thus retain a strict
known-urn preference under many
confidence inputs even after its
objective red chance is reduced.
No numerical confidence estimate
or empirical switching frequency
has been inferred from this
algebraic example.
\end{proof}
\end{annexitem}

\begin{annexitem}{annex:p17.ambiguity.other.criteria}{thm:p17-ambiguity-other-criteria}{Raw-state criteria and the complete regret counterexample}
\begin{proof}
Each of \(F_1,\ldots,F_4\) has
at least one utility-zero state
and at least one utility-one state.
Their state minima are therefore
all zero and maxima all one.
Any fixed weighted combination
of these extrema is identical
for all four acts.
For menu \(\{F_1,F_2\}\), the
statewise best vector is \((1,1,0)\);
its regret vectors are
\[
                     (0,1,0),\qquad(1,0,0).
\]
For menu \(\{F_3,F_4\}\), the
statewise best vector is \((1,1,1)\),
and the regret vectors are the
same two vectors. Both maximum
regrets equal one in each pair.
If all four acts are included
in one menu, the best vector is
still \((1,1,1)\), and every act
has some regret-one state, so
they remain tied by maximum
regret. These calculations prove
the exact claims about the
original urn comparisons.

For any fixed nonempty finite
menu, \({\rm Reg}_{\mathcal M}(x)\)
is a finite real number.
Ordering its negative gives a
complete transitive ranking.
In the stated two-state menu,
the statewise best payoff is
\((1,1)\). The respective maximum
regrets of
\[
             x=(1,1),\quad y=(0,1),\quad
             x'=(1,0),\quad y'=(0,0)
\]
are \(0,1,1,1\).
Thus \(x\succ y\) but \(x'\sim y'\).
The pair \(x',y'\) is obtained
by changing the common second
coordinate of \(x,y\) from one
to zero, so the sure-thing
condition would preserve strict
preference and is violated.

If the benchmark were recomputed
using only \(\{x',y'\}\), it
would instead be \((1,0)\), with
regrets zero and one. This
directly exhibits the distinction
between the pairwise and fixed-menu
conventions. The general statement
that all regret-based choices
preserve sure-thing or necessarily
lack a complete transitive ranking
would therefore be unwarranted
without those domain conventions.
\end{proof}
\end{annexitem}
